\documentclass[11pt]{amsart}
\usepackage[margin=1.1in]{geometry}
\usepackage{lmodern}
\usepackage{amsmath,amssymb,amsthm,mathtools}
\usepackage[mathscr]{eucal}
\usepackage{booktabs,enumitem,xcolor,array}
\usepackage[colorlinks=true,linkcolor=blue!50!black,citecolor=green!40!black,urlcolor=blue!60!black]{hyperref}
\hypersetup{
  pdftitle={Critical conductors of small archimedean data},
  pdfauthor={Nic Johns}
}
\setlength{\emergencystretch}{2.5em}

% ---------------------------------------------------------------- macros
\newcommand{\RR}{\mathbb R}
\newcommand{\ZZ}{\mathbb Z}
\newcommand{\CC}{\mathbb C}
\newcommand{\QQ}{\mathbb Q}
\newcommand{\dd}{\,\mathrm d}
\newcommand{\eps}{\varepsilon}
\renewcommand{\Re}{\operatorname{Re}}
\renewcommand{\Im}{\operatorname{Im}}
\newcommand{\supp}{\operatorname{supp}}
\newcommand{\one}{\mathbf 1}
\newcommand{\FT}[1]{\widehat{#1}}
\newcommand{\hF}{\widehat F}
\newcommand{\xin}[1]{\xi_{#1}}
\newcommand{\GR}{\Gamma_{\mathbb R}}
\newcommand{\GC}{\Gamma_{\mathbb C}}
\newcommand{\Oinf}{\Omega_\infty}
\newcommand{\Arch}{\mathcal A}
\newcommand{\Tc}{\mathcal T}
\newcommand{\Kad}{\mathcal K}
\newcommand{\Cone}{\mathcal C}
\newcommand{\ConeOPS}{\mathcal C_{\mathrm{OPS}}}
\newcommand{\kstar}{\kappa^{*}}
\newcommand{\kOPS}{\kappa^{*}_{\mathrm{OPS}}}
\newcommand{\qmin}{q_{\min}}
\newcommand{\Frep}{F_{\mathrm{rep}}}
\newcommand{\pz}{p_\zeta}

\newtheorem{theorem}{Theorem}[section]
\newtheorem{proposition}[theorem]{Proposition}
\newtheorem{lemma}[theorem]{Lemma}
\newtheorem{corollary}[theorem]{Corollary}
\newtheorem{fact}[theorem]{Fact}
\newtheorem{question}[theorem]{Question}
\theoremstyle{definition}
\newtheorem{definition}[theorem]{Definition}
\theoremstyle{remark}
\newtheorem{remark}[theorem]{Remark}
\newtheorem{observation}[theorem]{Numerical observation}
\numberwithin{equation}{section}

\title[Critical conductors]{Critical conductors of small archimedean data}
\author{Nic Johns}
\date{October 2026}
\subjclass[2020]{Primary 11R42; Secondary 11M26, 11R29, 90C05, 65G20}
\keywords{Weil explicit formula, conductor bounds, discriminant bounds, linear programming, computer-assisted proof}

\begin{document}

\begin{abstract}
For an archimedean type $\GR^a\GC^b$ with one simple pole, the least conductor at which the explicit formula admits a
solution by positive measures on both sides is the optimum of the explicit-formula method of Odlyzko, Poitou and Serre for
that type. We bracket it by computer-assisted certificates for $\GR^2$ and $\GC$, $4.98485094226764\le\qmin\le4.9994844$ and
$2.981236704765\le\qmin\le2.997174$, and compare the computed values for seven types with the smallest discriminants of
number fields with the same archimedean data. With the gap $\log2/2\pi$ common to all data, in every case tested the genuine
Dedekind zeta function lies strictly above its critical conductor, while for the data of the Riemann zeta function
$1-7.3\cdot10^{-906}\le\qmin$, and $\qmin\le1$ under the Riemann hypothesis. The two certified brackets hold for every gap up to
$\log4.049/2\pi$, resp.\ $\log3.011/2\pi$; so the data of $\QQ(\sqrt5)$ and $\QQ(\sqrt{-3})$ are not critical at their natural gaps
$\log4/2\pi$ and $\log3/2\pi$ either. We also certify an
admissible pair for the data of $\QQ(i)$, and that the pole weight of $\zeta$'s data, the residue of $-\zeta'/\zeta$ at $s=1$, is pinned:
no admissible pair exists for weights $r<1-5.5794\cdot10^{-40}$ or $r>1+7.8551\cdot10^{-21}$.
\end{abstract}

\maketitle

\noindent\textbf{Status.} This is a short computational note accompanying the companion paper \cite{PaperI}; it has not been
peer-reviewed. Its theorems and propositions are computer-assisted certificates, each with its script and parameter file
(Appendix~\ref{app:anc}), and are proved. The comparison of critical conductors with minimal discriminants
(Section~\ref{sec:observations}), the slack as a function of the gap (Numerical observation~\ref{obs:gap-profile}) and the evidence on
uniqueness at criticality are numerical observations: computed by finite
linear programmes, not certified, and not used in any proof. The comparisons use the common gap $\xin2$. For $\QQ(\sqrt5)$ and $\QQ(\sqrt{-3})$ the
certified conclusions also hold at their natural gaps (Corollary~\ref{cor:natural-gaps}); for the other types the computations
were made at the gap $\xin2$ only (Remark~\ref{rem:gap}). Whether the set of admissible pairs at a critical conductor is a single pair is open for
every type, $\zeta$'s included (Question~\ref{conj:Ggen}).

%======================================================================
\section{Introduction}\label{sec:intro}
%======================================================================

Explicit-formula bounds for discriminants and conductors \cite{Odl76,Poi77,DyD80,Odl90} test the explicit formula of an
$L$-function against a function $F\ge0$ whose Fourier transform is non-negative beyond the first prime. Miller \cite{Mil02} used
the same positivity, with carefully chosen test functions, to show that $L$-functions of many archimedean types have a lower first
zero than $\zeta$, or do not exist. For fixed archimedean data, the best bound
that the method can give is the least conductor $\qmin$ at which the explicit formula admits a solution by positive measures
on both sides, and by linear-programming duality it is $e^{-2\pi\kstar}$, where $\kstar$ is the slack of the data at conductor
$1$ \cite{PaperI}. For the data of $\zeta$ (the factor $\GR$ and one simple pole) the companion paper certifies
$\kstar\le1.1508391\cdot10^{-906}$, so that $\qmin\ge1-7.3\cdot10^{-906}$, while $\qmin\le1$ under the Riemann hypothesis; only the
first of these two inequalities is unconditional. It also shows that the data of $\QQ(\sqrt5)$ and $\QQ(\sqrt{-3})$
are not critical, with the same gap $\xin2$ and also at their natural gaps, which are larger because the prime-ideal norms of these
fields start at $4$ and $3$ (Fact~\ref{thm:notcritical}). Here we state the same certificates for the critical conductors at every gap
(Corollary~\ref{cor:natural-gaps}) and discuss the role of the gap (Remark~\ref{rem:gap}).

This note records the corresponding computations for other archimedean types $\GR^a\GC^b$ with one simple pole, in the spirit
of Odlyzko's problem of determining the optimal explicit-formula bounds \cite{Odl90}. We certify brackets for the critical
conductors of $\GR^2$ and $\GC$ (Theorem~\ref{thm:GR2} and Proposition~\ref{prop:GC}), valid for every gap up to slightly beyond the
natural gaps of $\QQ(\sqrt5)$ and $\QQ(\sqrt{-3})$ (Corollary~\ref{cor:natural-gaps}); the lower ends improve the
corresponding degree-$2$ entries of Odlyzko's tables (Remark~\ref{rem:odlyzko}). We certify an admissible pair for the data of
$\QQ(i)$ (Proposition~\ref{prop:Qi}) and a two-sided pin of the pole weight of $\zeta$'s data (Theorem~\ref{thm:polepin}).
For seven types, computed critical conductors are compared with the smallest discriminants of number fields with the same
archimedean data (Numerical observation~\ref{obs:critical-conductors}): with the gap $\xin2$, in every case the genuine field lies
strictly above its critical conductor, and the computed critical objects are not arithmetic.

%======================================================================
\section{Setting}\label{sec:setting}
%======================================================================

We recall the definitions and results of \cite{PaperI} that are used; statements quoted from it are called Facts. We write
$\hF(\xi)=\int_\RR F(t)e^{-2\pi i\xi t}\dd t$, $\xin n=\log n/(2\pi)$, and $\xin2=\log2/(2\pi)$ for the \emph{gap}, which is the same
for all archimedean data below, except in Corollary~\ref{cor:natural-gaps} (Remark~\ref{rem:gap}). The test class
$\Tc=\bigcup_{0<\delta<1/2}\Tc_\delta$ consists of the even functions $F$, real on $\RR$ and analytic in a neighbourhood of
$\{\lvert\Im z\rvert\le\frac12+\delta\}$ for some $\delta$, with $\sup(1+\lvert z\rvert)^2\lvert F(z)\rvert<\infty$ there.

For a degree $d\ge1$, shifts $\kappa=(\kappa_1,\dots,\kappa_d)$, a conductor $q>0$ and a pole weight $r>0$ put
\begin{gather*}
\Arch_{\kappa,q,r}(F):=r\bigl(F(\tfrac i2)+F(-\tfrac i2)\bigr)+\int_\RR F(t)\rho(t)\dd t,\\
\rho(t):=\frac1{2\pi}\Bigl[\log q+\sum_{j\le d}\bigl(\Re\psi\bigl(\tfrac{\frac12+\kappa_j+it}{2}\bigr)-\log\pi\bigr)\Bigr],
\end{gather*}
with $\psi=\Gamma'/\Gamma$. The data of $\zeta$ are $d=1$, $\kappa=(0)$, $q=r=1$; $\GR^2$ has $\kappa=(0,0)$; $\GC(s)\propto\GR(s)\GR(s+1)$ has
$\kappa=(0,1)$; in general $\GR^a\GC^b$ has $a+b$ shifts $0$ and $b$ shifts $1$. The \emph{pole weight} $r$ is the residue at $s=1$ of
$-L'/L$ for an $L$-function $L$ with these data, that is, the coefficient of $F(\frac i2)+F(-\frac i2)$ in its explicit formula; it
is not the residue of $L$ itself, which changes when $L$ is multiplied by a non-zero constant, while $-L'/L$ and the explicit
formula do not. A simple pole gives $r=1$. An \emph{admissible pair} for $\Arch'=\Arch_{\kappa,q,r}$ is a
pair $(\mu,\nu)$ of positive Radon measures, $\mu$ even on $\RR$ and $\nu$ on $[\xin2,\infty)$, such that
$\Arch'(F)=\int F\dd\mu+\frac1\pi\int\hF\dd\nu$ for every $F\in\Tc$, with absolutely convergent integrals; $\Kad(\Arch')$ is the set of
admissible pairs. Put
\[
\Cone:=\{F\in\Tc:F\ge0\text{ on }\RR,\ \hF\ge0\text{ on }[\xin2,\infty)\},\qquad
\ConeOPS:=\{F\in\Tc:F\ge0\text{ and }\hF\ge0\text{ on }\RR\},
\]
$\kstar(\Arch'):=\inf\{\Arch'(F):F\in\Cone,\ \int F=1\}$, and $\kOPS(\Arch')$ likewise with $\ConeOPS$. Since
$\Arch_{\kappa,q',r}=\Arch_{\kappa,q,r}+\frac{\log(q'/q)}{2\pi}\int$, feasibility is monotone in $q$, and
$\kstar(\Arch_{\kappa,q,r})=\kstar(\Arch_{\kappa,1,r})+\frac{\log q}{2\pi}$. The \emph{critical conductor} is
$\qmin(\kappa,r):=\inf\{q:\Kad(\Arch_{\kappa,q,r})\ne\emptyset\}$.

\emph{Other gaps.} For $g>0$, replacing $[\xin2,\infty)$ by $[g,\infty)$ in these definitions defines admissible pairs \emph{at the gap $g$},
the cone $\Cone_g$, the slack $\kstar_g$ and the critical conductor $q_{\min,g}$; for $g=\xin2$ they are the objects above. A larger gap
allows fewer pairs and gives a larger cone, so $q_{\min,g}$ is non-decreasing and $\kstar_g$ is non-increasing in $g$.

\begin{lemma}[Weak duality and the critical conductor]\label{lem:weak-duality}
\textup{(a)} If $(\mu,\nu)$ is admissible for $\Arch'$ and $F\in\Cone$, then for Borel sets $I\subseteq\RR$ and
$J\subseteq[\xin2,\infty)$, $(\inf_IF)\mu(I)+\frac1\pi(\inf_J\hF)\nu(J)\le\Arch'(F)$; in particular $\Arch'(F)\ge0$, and if $\int F=1$ no
admissible $\mu$ satisfies $\mu\ge\lambda\dd t$ with $\lambda>\Arch'(F)$. \textup{(b)} An admissible pair exists if and only if
$\kstar(\Arch')\ge0$; hence $\qmin(\kappa,r)=e^{-2\pi\kstar(\Arch_{\kappa,1,r})}$.
\end{lemma}

\begin{proof}
For $\zeta$'s data, (a) is \cite[Lemma~2.5]{PaperI} and (b) is \cite[Theorem~2.7 and Proposition~2.10]{PaperI}. Part (a) needs one line:
for an admissible pair and $F\in\Cone$, both integrands in $\Arch'(F)=\int F\dd\mu+\frac1\pi\int\hF\dd\nu$ are non-negative where the
measures live. For (b), the proofs in \cite{PaperI} use only that the archimedean functional is $\lVert\cdot\rVert_\delta$-continuous on each
$\Tc_\delta$ and the bounds for its values on the test functions of the proof of \cite[Lemma~2.6]{PaperI}, which follow from the growth of
the archimedean weight. Both hold for $\Arch_{\kappa,q,r}$: its pole term $r(F(\frac i2)+F(-\frac i2))$ is $\lVert\cdot\rVert_\delta$-continuous, and
$\rho$ is continuous with $\rho(t)=\frac d{2\pi}\log\lvert t\rvert+O(1)$, as for $\zeta$ with $d=1$. So these proofs apply verbatim; the
conductor enters as $\frac{\log q}{2\pi}$ times Lebesgue measure, which gives the formula for $\qmin$.
\end{proof}

\begin{lemma}[Herglotz form of an admissible pair]\label{lem:herglotz}
Let $\tilde\nu\ge0$ be a measure on $[2,\infty)$ such that $\sigma:=\tilde\nu-r\one_{[2,\infty)}(x)\dd x$ satisfies
$\int x^{-1/2}\dd\lvert\sigma\rvert(x)<\infty$, and put $\mathcal L(s):=\int x^{-s}\dd\tilde\nu(x)$ for $\Re s>1$. Then $\mathcal L$ extends
continuously to $\Re s\ge\frac12$, $s\ne1$, as $r2^{1-s}/(s-1)+\int x^{-s}\dd\sigma$. Let $\mu(t):=\rho(t)-\frac1\pi\Re\mathcal L(\frac12+it)$,
and let $\nu$ be the image of $x^{-1/2}\tilde\nu$ under $x\mapsto\log x/2\pi$. Then $\Arch_{\kappa,q,r}(F)=\int F\mu\dd t+\frac1\pi\int\hF\dd\nu$
for every $F\in\Tc$. Hence $(\mu\dd t,\nu)$ is admissible for $\Arch_{\kappa,q,r}$ if and only if $\mu\ge0$.
\end{lemma}

\begin{proof}
For $r=1$ and the data $\GR$, $\GR^2$, $\GC$ this is \cite[Lemma~5.7]{PaperI}. Its proof shifts the contour in
$\frac1\pi\int\hF\dd\nu=\frac1\pi\int F(w)\mathcal L(\frac12+iw)\dd w$ past the simple pole of $\mathcal L$ at $s=1$, whose residue, here $r$,
produces the pole term $r(F(\frac i2)+F(-\frac i2))$, while the archimedean weight enters only through the definition of $\mu$. So the proof
applies verbatim for every $r>0$ and every $\kappa$.
\end{proof}

\emph{Explicit pairs.} The pairs certified below have the following form: on $x\ge1$,
\[
\tilde\nu=\one_{[2,\infty)}(x)\dd x-\sum_{i<n}f_i\one_{[e^{v_i},e^{v_{i+1}}]}(x)\dd x+\sum_b(\text{bump}_b)-\sum_ma_mx^{-m}\one_{(X,\infty)}(x)\dd x,
\]
with $0\le f_i\le1$, cell edges $\log2=v_0<\dots<v_n$, $X=e^{v_n}$, $a_m\ge0$, and bumps
$w_bx^{1/2}e^{-u_b/2}\phi(\log x-u_b)\dd(\log x)$ with $w_b\ge0$, where $\phi$ is the density of a rescaled Irwin--Hall
distribution, supported on $[-\eps,\eps]$, with Fourier transform $S(t)=\operatorname{sinc}(\eps t/k)^k$ (for $\eps=0$, atoms). Then
$\tilde\nu\ge0$ with support in $[2,\infty)$ provided $u_b-\eps\ge\log2$ for every bump and $1-\sum_ma_mx^{-m}\ge0$ on $(X,\infty)$; since
$a_m\ge0$ the last function is increasing, so $x=X$ suffices. In Lemma~\ref{lem:herglotz} (with $r=1$), $\mu$ is explicit:
\[
\mu(t)=\rho(t)+\frac1\pi\sum_{j=0}^nC_je^{v_j/2}E(t,v_j)-\frac{S(t)}\pi\sum_bm_b\cos(tu_b)+\sum_m\frac{a_m}\pi X^{1/2-m}T_m(t),
\]
with $E(t,v)=\Re[e^{itv}/(\frac12+it)]$, $T_m(t)=\Re[e^{-itL}/(m-\frac12+it)]$, $L=\log X$, $C_0=1-f_0$, $C_j=f_{j-1}-f_j$
($0<j<n$), $C_n=f_{n-1}$ and $m_b=w_be^{-u_b/2}$. The verifier proves $\mu\ge0$ on a finite interval $[0,T]$ by mean-value steps,
$\mu(t)\ge\mu(t_0)-(h+d)\sup\lvert\mu'\rvert$ on $[t_0-d,t_0+h+d]$, with $\mu'$ enclosed on the whole step in ball arithmetic. Beyond
$T$ it uses monotone majorants: $\rho$ is increasing, while $\lvert E(t,v)\rvert\le(\frac14+t^2)^{-1/2}$,
$\lvert S(t)\rvert\le\min(1,(k/\eps t)^k)$ and $\lvert T_m(t)\rvert\le((m-\frac12)^2+t^2)^{-1/2}$ are decreasing, so for $t\ge T$, $\mu(t)$ is at
least $\rho(T)$ minus the sum of these majorants at $T$, with the coefficients in absolute value.

\emph{Gaussian--Laguerre certificates.} Let $f_k(u):=L_k^{(-1/2)}(2\pi u^2)e^{-\pi u^2}$; these are even Hermite functions, with
$\widehat{f_k}=(-1)^kf_k$, the basis with prescribed double roots used for sign-uncertainty problems in \cite{CG19}. The dual
certificates below are functions
\[
\Frep(t):=\sum_{k\le K}a_kf_k(t/s)+\eps\bigl(e^{-\pi t^2/(16s^2)}+e^{-16\pi t^2/s^2}\bigr)
\]
with explicit rationals $a_k$. Writing $F:=\sum_ka_kf_k(t/s)$, $X=2\pi t^2/s^2$ and $Y=2\pi s^2\xi^2$, one has $F=e^{-X/2}P(X)$ and
$\hF=s\,e^{-Y/2}Q(Y)$ with $P,Q\in\QQ[X]$, and exact Sturm sequences decide the positivity of $P$ on $(0,\infty)$ and of $Q$ on the
range of $Y$ in question. The two Gaussians of the repair term are positive with positive Fourier transforms, so $\Frep\ge F$
and $\FT F_{\mathrm{rep}}\ge\hF$ pointwise; the repair term is not needed for these positivity statements, and it adds the explicit
lower bounds $\Frep(t)\ge\eps e^{-\pi t^2/(16s^2)}$ and, where $\hF\ge0$, $\FT F_{\mathrm{rep}}(\xi)\ge\frac{\eps s}4e^{-\pi s^2\xi^2/16}$.
$\Arch'(\Frep)$ is one rigorous ball-arithmetic integral of $\Frep$ against the digamma weight on a finite interval, plus an
explicit tail bound.

\begin{fact}[Results of \cite{PaperI} on $\zeta$, $\QQ(\sqrt5)$ and $\QQ(\sqrt{-3})$]\label{thm:notcritical}
\textup{(a)} For $\zeta$'s data, $-2.7112\cdot10^{-3}\le\kstar\le1.1508391\cdot10^{-906}$ \textup{(}\cite[Theorem~C]{PaperI}\textup{)}; and
$\kstar\le\kOPS\le9.9462\cdot10^{-41}$, certified by a Gaussian--Laguerre certificate with $K=128$, $s=13$, $\eps=10^{-56}$ whose
transform is positive on all of $\RR$ \textup{(}\cite[Proposition~5.3]{PaperI}\textup{)}. \textup{(b)} For the archimedean data of $\QQ(\sqrt5)$
\textup{(}$\GR^2$, $q=5$, $r=1$\textup{)} and every gap $g\in(0,\xin{4.04915}]$, in particular for $\xin2$ and for the natural gap $\xin4$ of the
field, $1.64\cdot10^{-5}\le\kstar_g\le4.83\cdot10^{-4}$; for the data of $\QQ(\sqrt{-3})$ \textup{(}$\GC$, $q=3$, $r=1$\textup{)} and every gap
$g\in(0,\xin{3.011664}]$, in particular for $\xin2$ and for the natural gap $\xin3$, $1.50\cdot10^{-4}\le\kstar_g\le9.99\cdot10^{-4}$. In both cases
some admissible pair has zero measure bounded below by a positive multiple of $\dd t$, and the set of admissible pairs is not a
singleton \textup{(}\cite[Proposition~5.9]{PaperI}\textup{)}.
\end{fact}

The proof of (b) in \cite{PaperI} combines the lower bound of Theorem~\ref{thm:GR2} below (resp.\ Proposition~\ref{prop:GC}), read at
$q=5$ (resp.\ $q=3$), with the explicit pair of the same statement. Corollary~\ref{cor:natural-gaps} below states the same certificates
in terms of the critical conductors $q_{\min,g}$.

%======================================================================
\section{Brackets for two critical conductors}\label{sec:brackets}
%======================================================================

\begin{theorem}[A non-arithmetic critical point, bracketed]\label{thm:GR2}
For the data $\GR^2$ with one simple pole \textup{(}$\kappa=(0,0)$, $r=1$\textup{)},
\[
4.98485094226764\le\qmin\le4.9994844 .
\]
More precisely: an explicit Gaussian--Laguerre certificate $\Frep\in\Cone$ with $K=80$ and $s=7$ has
$\Arch_{(0,0),1,1}(\Frep)/\int\Frep\le-0.25566705789831547$, so no admissible pair exists for $q<4.98485094226764$; and an explicit
pair of the form above at $q_0=e^{1.6093347792651136}\le4.9994844$ has $\mu\ge6.59\cdot10^{-9}$ on $[0,800]$ and $\mu\ge1.79$ on
$[800,\infty)$.
\end{theorem}

\begin{proof}[Proof \textup{(computer-assisted)}]
Lower bound: with the notation above, exact Sturm sequences show that $P$ has no root in $(0,\infty)$, that $Q$ has no root
in $(Y_{\mathrm{lo}},\infty)$ for a rational $Y_{\mathrm{lo}}$ below the image of $\xin2$, and that the relevant end values and leading
coefficients are positive; so $\Frep>0$ on $\RR$ and $\FT F_{\mathrm{rep}}>0$ on $[\xin2,\infty)$, and $\Frep\in\Cone$. In fact $Q$ has no
root in $(0,\infty)$ and $Q(0)>0$, so $\FT F_{\mathrm{rep}}>0$ on all of $\RR$ and $\Frep\in\ConeOPS$. The functional is one
Arb integral at $448$ bits (\texttt{python} \nolinkurl{general/verify_general.py} \nolinkurl{general/params/gammaR2_K80_s7.json}). Then
$\Arch_{(0,0),q,1}(\Frep)<0$ for $\log q<2\pi\cdot0.25566705789831547$, and Lemma~\ref{lem:weak-duality} applies. Upper bound:
\texttt{python} \nolinkurl{general/verify_pair.py} \nolinkurl{general/params/pair_gammaR2_X240.json} checks $\tilde\nu\ge0$ and
$\supp\tilde\nu\subseteq[2,\infty)$ (the first edge is exactly $\log2$ in the $\log x$ variable), and proves $\mu>0$ on $[0,800]$ by
$17\,214$ mean-value steps with $\mu'$ enclosed in Arb ($128$ bits) on each whole step, and on $[800,\infty)$ by the monotone
majorants described in Section~\ref{sec:setting}. Lemma~\ref{lem:herglotz} gives admissibility at $q_0$, and feasibility for $q\ge q_0$
follows by monotonicity.
\end{proof}

\begin{proposition}[The critical conductor of $\GC$]\label{prop:GC}
For the data $\GC$ with one simple pole \textup{(}$\kappa=(0,1)$, $r=1$\textup{)}, $2.981236704765\le\qmin\le2.997174$.
\end{proposition}

\begin{proof}[Proof \textup{(computer-assisted)}]
As for Theorem~\ref{thm:GR2}, with a Gaussian--Laguerre certificate for $\kappa=(0,1)$ ($K=64$, functions $f_k(t/9)$;
\nolinkurl{general/params/gammaC_K64_s9.json}, which prints $\qmin\ge2.9812367047650810559$) and an explicit pair of the form above
at $\log q_0=1.0976698108720329$ (\nolinkurl{general/params/pair_gammaC_X40.json}), whose zero density satisfies the certified floor
$\mu\ge6.67\cdot10^{-8}$ on $\RR$. As in Theorem~\ref{thm:GR2}, the transform of the dual certificate is positive on all of $\RR$.
\end{proof}

Neither bracket uses the value of the gap in an essential way: the dual certificates have transforms positive on all of $\RR$,
and the prime measures of the explicit pairs start well beyond $2$.

\begin{corollary}[The brackets at other gaps]\label{cor:natural-gaps}
\textup{(a)} For the data $\GR^2$ with one simple pole, $4.98485094226764\le q_{\min,g}\le4.9994844$ for every gap $g$ with
$0<g\le\xin{4.04915}$; and the data of $\QQ(\sqrt5)$ \textup{(}$\GR^2$, $q=5$, $r=1$\textup{)} have $1.64\cdot10^{-5}\le\kstar_g\le4.83\cdot10^{-4}$
for these $g$.
\textup{(b)} For the data $\GC$ with one simple pole, $2.981236704765\le q_{\min,g}\le2.997174$ for every gap $g$ with
$0<g\le\xin{3.01166}$; and the data of $\QQ(\sqrt{-3})$ \textup{(}$\GC$, $q=3$, $r=1$\textup{)} have $1.50\cdot10^{-4}\le\kstar_g\le9.99\cdot10^{-4}$
for these $g$.

\noindent In particular, at the natural gaps $\xin4$ of $\QQ(\sqrt5)$ and $\xin3$ of $\QQ(\sqrt{-3})$, some admissible pair has zero measure
bounded below by a positive multiple of $\dd t$, and these data are not critical.
\end{corollary}

\begin{proof}[Proof \textup{(computer-assisted)}]
At a gap $g$, weak duality takes the form of Lemma~\ref{lem:weak-duality}(a) with $\Cone_g$ in place of $\Cone$: for an admissible
pair at the gap $g$ and $F\in\Cone_g$, both integrands in $\Arch'(F)=\int F\dd\mu+\frac1\pi\int\hF\dd\nu$ are non-negative on the supports
of the measures. The identity of Lemma~\ref{lem:herglotz} does not involve the gap, so its pair is admissible at the gap $g$ if and only
if $\mu\ge0$ and $\tilde\nu$ is carried by $[e^{2\pi g},\infty)$.

(a) The dual certificate of Theorem~\ref{thm:GR2} lies in $\ConeOPS\subseteq\Cone_g$ (its log reports
\texttt{F\_rep\^{} > 0 on R (OPS cone): True}), so no admissible pair exists at any gap for $q<4.98485094226764$, and at $q=5$ it
gives $\kstar_g\le-0.25566705789831547+\log5/2\pi\le4.83\cdot10^{-4}$. Run with the option \nolinkurl{--gap-x} \nolinkurl{4}, the verifier
of the explicit pair of Theorem~\ref{thm:GR2} proves, by exact comparisons of the cell edges and outward-rounded bounds for the bump
supports, that \texttt{supp nu\textasciitilde{} in [x0, oo) with x0 >= 4.049150}, and reports \texttt{admissible at gap xi\_4 (prime
measure on [4, oo)): True}. So the pair is admissible at every gap $g\le\xin{x_0}$, and $q_{\min,g}\le4.9994844$ there. With $\mu$
replaced by $\mu+\frac1{2\pi}\log(5/q_0)$ the same $\tilde\nu$ gives an admissible pair for the data of $\QQ(\sqrt5)$ whose zero measure
is at least $1.64\cdot10^{-5}\dd t$, so $\kstar_g\ge1.64\cdot10^{-5}$. The script \nolinkurl{general/corollary_slack.py} combines the two
runs: \texttt{python} \nolinkurl{general/corollary_slack.py} \nolinkurl{general/params/corollary_qsqrt5.json} \nolinkurl{--gap-x}
\nolinkurl{4} prints $4.984850942267643\le\qmin\le4.999484361$ and $1.6420748\cdot10^{-5}\le\kstar\le0.00048294147$, and that both hold
for every gap $g$ with $e^{2\pi g}\le x_0$.

(b) The same, with Proposition~\ref{prop:GC}: its dual certificate lies in $\ConeOPS$; the verifier of its pair, run with
\nolinkurl{--gap-x} \nolinkurl{3}, proves $x_0\ge3.011664$ and admissibility at the gap $\xin3$; and the run of
\nolinkurl{general/corollary_slack.py} with \nolinkurl{general/params/corollary_qsqrtm3.json} \nolinkurl{--gap-x} \nolinkurl{3} prints
$0.00015006672\le\kstar\le0.00099854968$ for the data of $\QQ(\sqrt{-3})$.
\end{proof}

\begin{observation}[The slack as a function of the gap]\label{obs:gap-profile}
The slack $\kstar_g$ at $g=\xin n$ was also computed by exact sums-of-squares linear programmes in the Gaussian--Laguerre basis
($K\le80$, an interior-point solver at $640$ bits, duality gaps below $10^{-24}$; every solution with $K\ge32$ was re-checked at
$256$ bits on fine grids, with refined local minima); these computations are not part of the ancillary files. For $\zeta$'s data
and for those of $\QQ(\sqrt5)$ and $\QQ(\sqrt{-3})$, $\kstar_g$ is constant for $0<g\le\xin{n_1}$, where $n_1$ is the first prime atom
of the extremal pair of the linear programme, and equals there the slack $\kOPS$ of the cone $\ConeOPS$; beyond $\xin{n_1}$ it
decreases roughly linearly in $g$, with slope between $-1.85$ and $-1.5$, and becomes negative. For $\zeta$'s data the computed values of the
constant decay like $10^{-0.3K}$, consistent with $0$ (unconditionally only $-2.7112\cdot10^{-3}\le\kstar\le1.1508391\cdot10^{-906}$ is
known, and $\kstar\ge0$ holds if admissible pairs exist, for instance under the Riemann hypothesis), and $n_1=2$; the slack is $-0.0470$ at $\xin{2.5}$ and $-0.0793$ at $\xin3$. For the
data of $\QQ(\sqrt5)$ the constant is $4.8294\cdot10^{-4}$ and $n_1\approx4.3201$; the slack vanishes near $n\approx4.328$ and is
$-0.0372$ at $\xin5$. For the data of $\QQ(\sqrt{-3})$ the constant is $9.9855\cdot10^{-4}$ and $n_1\approx3.2442$; the slack vanishes
near $n\approx3.255$ and is $-0.0570$ at $\xin4$. Where Corollary~\ref{cor:natural-gaps} applies, every computed value with $K\ge32$ lies
in its bracket. In these examples the extremal pair for $\zeta$'s data is arithmetic, its first prime atom being the prime $2$, while for
the two fields the first atom falls between prime-ideal norms and the constant is positive; we draw no general conclusion from this.
\end{observation}

\begin{remark}[Comparison with Odlyzko's tables]\label{rem:odlyzko}
If the generalised Riemann hypothesis holds for the Dedekind zeta function of a real (resp.\ imaginary) quadratic field $k$,
then its own pair of measures (zeros on the zero side, prime ideals and their powers on the prime side, which lies beyond
$\xin2$ because every norm is at least $2$) is admissible for the data $\GR^2$ (resp.\ $\GC$) at $q=\lvert d_k\rvert$, $r=1$, by the
explicit formula for the Dedekind zeta function of $k$, which holds on $\Tc$ by the argument used for $\zeta$; so
$\lvert d_k\rvert\ge\qmin$. In the normalisation of root discriminants, the lower ends of Theorem~\ref{thm:GR2} and
Proposition~\ref{prop:GC} are $4.98485^{1/2}=2.2327$ and $2.98123^{1/2}=1.7266$, against the GRH bounds $2.225$ and $1.722$ in
degree $2$ of \cite[Tables~4 and~3, p.~134]{Odl90} (see also \cite{Odl76,DyD80}); the minimal root discriminants are
$5^{1/2}=2.236$ and $3^{1/2}=1.732$. Squared, the tabulated bounds are $2.225^2\approx4.95<5$ and $1.722^2\approx2.97<3$, consistent with our
brackets. In the sense of Cohn and Triantafillou \cite{CT22}, the explicit pairs of Theorem~\ref{thm:GR2} and
Proposition~\ref{prop:GC} are dual certificates of non-sharpness: no test function in $\Cone$ can prove a conductor bound above
$4.9994844$, resp.\ $2.997174$, so for these data the method cannot reach the minimal discriminants $5$ and $3$.
\end{remark}

\begin{remark}[The role of the gap]\label{rem:gap}
Except in Corollary~\ref{cor:natural-gaps}, this note uses the gap $\xin2=\log2/2\pi$, the first prime of $\QQ$, for every archimedean
type. For the data of a number field $k\ne\QQ$ this gap need not be natural. In $\QQ(\sqrt5)$ the primes $2$ and $3$ are inert, so the
norms of the prime ideals are $4,5,9,11,11,\dots$; in $\QQ(\sqrt{-3})$ the first norm is $3$. At the gap $\xin4$, resp.\ $\xin3$,
positivity of $\hF$ is required only beyond that point, so the cone is larger and $\kstar$ can only decrease; and if the generalised
Riemann hypothesis holds for $\zeta_k$, the field's own pair is admissible at that gap, so there $\kstar\ge0$. So positive slack at
the gap $\xin2$ does not by itself give positive slack at the natural gap. For these two fields it does
(Corollary~\ref{cor:natural-gaps}): the dual certificates have transforms positive on all of $\RR$, and the explicit pairs put no prime
mass below $x=4.049$, resp.\ $x=3.011$. The margins are small: the critical conductor of the data $\GR^2$ with one simple pole, which
are those of $\QQ(\sqrt5)$ up to the conductor, lies in $[4.9848,4.9995]$ at every gap up to $\xin{4.049}$, so its distance below the
discriminant $5$ is between $10^{-4}$ and about $3\cdot10^{-3}$ in relative terms. For the other types of Numerical
observation~\ref{obs:critical-conductors} the computations were made at the gap $\xin2$ only, and their comparison with minimal
discriminants may depend on the gap. In general the critical conductor can depend on the gap, and so can Question~\ref{conj:Ggen}.
\end{remark}

\begin{proposition}[An admissible pair for the data of $\QQ(i)$]\label{prop:Qi}
For the data of $\QQ(i)$ \textup{(}$\GC$, $q=4$, $r=1$\textup{)} there is an admissible pair of the form above with $\mu\ge0.0217$ on
$[0,300]$ and $\mu\ge0.4696$ on $[300,\infty)$. Hence $\kstar\ge0.0217$ for these data, and the set of admissible pairs is not a
singleton.
\end{proposition}

\begin{proof}[Proof \textup{(computer-assisted)}]
As for the upper bound in Theorem~\ref{thm:GR2}: the script \nolinkurl{general/verify_pair.py} with
\nolinkurl{general/params/pair_Qi_X10.json} checks $\tilde\nu\ge0$ with support in $[2,\infty)$ and proves the bounds for $\mu$ by
$1277$ mean-value steps with $\mu'$ enclosed in Arb, and by monotone majorants beyond $300$; the printed floor is
$\mu\ge0.0222324$ on $\RR$. Lemma~\ref{lem:herglotz} gives admissibility, and $\mu\ge0.0217\dd t$ gives $\kstar\ge0.0217$ by
Lemma~\ref{lem:weak-duality}. Adding to $\tilde\nu$ a small atom $w\delta_{x_1}$, $x_1\ge2$, changes $\mu$ by
$-(w/\pi)x_1^{-1/2}\cos(t\log x_1)$; for small $w>0$ the result is a second admissible pair.
\end{proof}

%======================================================================
\section{The pole weight of \texorpdfstring{$\zeta$}{zeta}'s data}\label{sec:pole}
%======================================================================

\begin{theorem}[The pole weight is pinned]\label{thm:polepin}
For $\zeta$'s data with pole weight $r$, that is $\Arch_r(F):=r(F(\frac i2)+F(-\frac i2))+\int F\Oinf$ with
$\Oinf(t)=\frac1{2\pi}(\Re\psi(\frac14+\frac{it}2)-\log\pi)$, there is no admissible pair if $r<1-5.5794\cdot10^{-40}$ or
$r>1+7.8550942\cdot10^{-21}$. Under the Riemann hypothesis, $r=1$ is feasible.
\end{theorem}

\begin{proof}[Proof \textup{(computer-assisted)}]
$\Arch_r(F)=\Arch_1(F)+(r-1)D$ with $D=F(\frac i2)+F(-\frac i2)$, so one certificate $F\in\Cone$ per side suffices. For the left
side, the $K=128$ certificate of Fact~\ref{thm:notcritical}(a) has $D>0$, and its run (\texttt{python}
\nolinkurl{general/verify_general.py} \nolinkurl{general/params/zeta_K128_s13.json}, also in the ancillary files of \cite{PaperI}) prints
$\Arch_1(F)/D\le5.57930810147\cdot10^{-40}$. For the right side, a certificate with $K=80$, $s=12$, $\eps=10^{-30}$, whose transform is
positive on $[\xin2,\infty)$ but negative on part of $(0,\xin2)$, has $D<0$ and $\Arch_1(F)/\lvert D\rvert\le7.85509419726\cdot10^{-21}$
(\texttt{python} \nolinkurl{general/verify_general.py} \nolinkurl{general/params/zeta_K80_s12_edge.json}). In both cases $\Arch_r(F)<0$ beyond the
stated thresholds, and Lemma~\ref{lem:weak-duality} applies. Under the Riemann hypothesis, $\zeta$'s own pair $\pz$ is admissible for
$r=1$ \cite{PaperI}.
\end{proof}

The same verifier certifies $\kOPS\le7.24651\cdot10^{-32}$ with $K=104$ (parameter file
\nolinkurl{general/params/zeta_K104_s13.json}); this certificate also has $D>0$, with $\Arch_1(F)/D\le4.06492480958\cdot10^{-31}$, which
gives the weaker left pin $r<1-4.0649249\cdot10^{-31}$. Like the bound $\kOPS\le9.9462\cdot10^{-41}$ for $K=128$, the bound for
$\kOPS$ needs no gap at $\xin2$ (the transform of the certificate is positive on all of $\RR$), so it is directly comparable with
the classical tables \cite[Table~4]{Odl90}. Only this pinning is certified; we make no claim about the shape of
the feasible region in other directions.

%======================================================================
\section{Numerical observations}\label{sec:observations}
%======================================================================

The observations below were computed but not certified, except where a certified statement is cited; none is used in a proof.

\begin{observation}[Critical conductors of other Gamma types]\label{obs:critical-conductors}
For each Gamma type with one simple pole, the critical conductor $e^{-2\pi\kstar_0}$, where $\kstar_0$ is the slack at
$q=1$, was computed by linear programming at finite $K$ (converged in $K$, and reproduced to $8$ digits of $\kstar_0$ by
an independent implementation; not certified except as in Theorem~\ref{thm:GR2} and Proposition~\ref{prop:GC}). We list it
against the smallest absolute discriminant of a number field with these archimedean data:
$\GC$: $2.981$ / $3$; $\GR^2$: $4.985$ / $5$; $\GR\GC$: $22.6$ / $23$; $\GR^3$: $48.2$ / $49$; $\GC^2$: $114.7$ / $117$;
$\GR^2\GC$: $267.8$ / $275$; $\GR^4$: $694.6$ / $725$. With the gap $\xin2$, in every case tested the genuine Dedekind zeta
function lies strictly above its critical conductor. For $\GR^2$ and $\GC$ this also holds at the natural gaps of $\QQ(\sqrt5)$ and
$\QQ(\sqrt{-3})$, and is certified (Corollary~\ref{cor:natural-gaps}); for the other five types the computations were made at the gap
$\xin2$ only, and the comparison may depend on the gap (Remark~\ref{rem:gap}). For $\zeta$, by contrast, $1-\qmin\le7.3\cdot10^{-906}$ (Fact~\ref{thm:notcritical}(a)),
and $\qmin\le1$ under the Riemann hypothesis.
\end{observation}

\begin{observation}[The near-critical objects are not arithmetic]\label{obs:nonarith}
Read from the parameter files: below $x\approx17.8$ the prime measure of the explicit pair of Theorem~\ref{thm:GR2} is a sum of
narrow bumps (half-width $0.05$ in $\log x$), the first ones centred near $x\approx4.257$, $4.277$, $8.410$, $8.591$, $11.57$ and $12.47$;
and the prime measure of the pair of Proposition~\ref{prop:Qi} has atoms at $x\approx3.089$, $3.113$, $4.693$, $4.731$, $5.425$ and
$5.469$, and so on. The computed critical object for the data $\GR^2$ (from the linear programmes of Numerical
observation~\ref{obs:gap-profile}) has zeros near $6.75$, $10.44$, $12.83$ and prime atoms near $x\approx4.3201$, $9.215$, $14.41$, and
so on. None of these positions is an integer, so none of these measures resembles the
prime-ideal measure of a number field (respectively of $\ZZ[i]$). So the critical object for these data appears to be
non-arithmetic; this reading is heuristic.
\end{observation}

\begin{observation}[Window masses near the critical point of $\GR^2$]\label{obs:windows}
For the data $\GR^2$, the largest mass that an admissible pair can put in four windows, two on the prime side (the images of
$x\in(2.05,3.95)$ and of $x\in(5,8.5)$) and two on the zero side ($(0.5,6)$ and $(7.6,9.7)$), was computed by finite linear
programmes ($K=48$) at conductors $\log q=\log\qmin+\delta$. It decreases as $\delta\downarrow0$, from $0.314$, $1.26$,
$0.078$ and $0.67$ at $\delta=10^{-3}$ to $0.0044$, $0.148$, $0.00085$ and $0.022$ at $\delta=10^{-5}$, while for the slack data of
$\QQ(\sqrt5)$ the same windows carry up to $0.60$, $1.57$, $0.20$ and $0.94$. These computations are not reproduced in the
ancillary files. They are consistent with a single admissible pair at the critical conductor, but these computations at $\delta>0$
do not show this: that would need a convergence statement for the dual problems as $\delta\to0$, which has only been computed.
Similarly, in three one-parameter families of Gamma factors, finite linear programmes ($K=48$) showed no visible fattening of the
set of admissible pairs at the critical conductor; finite computations of this kind can suggest, but not show, that a critical
fibre is a single pair.
\end{observation}

%======================================================================
\section{A question}\label{sec:question}
%======================================================================

For other archimedean data the question of uniqueness at criticality is open even numerically.

\begin{question}[Uniqueness at criticality for other Gamma factors]\label{conj:Ggen}
Let $\gamma=\GR^a\GC^b$ be a Gamma factor, with one simple pole at $s=1$ and the gap $\xin2$, and let $q_c(\gamma)$ be its critical
conductor, at which the slack vanishes. Is the set of admissible pairs at $q=q_c(\gamma)$ a single pair?
\end{question}

The critical conductor, and with it the question, can depend on the gap (Remark~\ref{rem:gap}). For $\zeta$'s own type the
question concerns Conjecture~U of \cite{PaperI} only when $\kstar=0$, for instance under the Riemann hypothesis and Conjecture~S of
\cite{PaperI}. Then $q_c=1$, admissible pairs exist, and the question asks whether there is only one; under the Riemann hypothesis,
which makes $\zeta$'s own pair admissible, this is Conjecture~U. If $\kstar\ne0$, it concerns a conductor $q_c\ne1$ and is unrelated
to Conjecture~U. For $\GR^2$ the critical conductor is bracketed by Theorem~\ref{thm:GR2}. A first step would be to certify, at the critical
conductor itself or at certified conductors converging to it, upper bounds for the window masses of Numerical
observation~\ref{obs:windows}.

\section*{AI-use statement}
This note was produced with substantial use of AI. The mathematical arguments, the numerical computations and the text were generated by instances of Claude Opus 5.5
(Anthropic) working as coordinated agents: one set of agents developed the arguments, separate agents acted as
independent hostile referees of each component, and the referees' requested repairs were incorporated. The numerical
certificates were computed by AI-written scripts that are included as supplementary material and can be rerun. The
author initiated and directed this project, designed its verification process, and reviewed the results and the
verification record, but did not independently check every step of the written proofs. The author takes full
responsibility for the content of this note, including any errors, and welcomes independent verification and
corrections.

\appendix
%======================================================================
\section{Ancillary files}\label{app:anc}
%======================================================================

Every computer-assisted statement names its script and parameter file; the scripts need Python $\ge3.10$ with python-flint
$0.9.0$ and mpmath $1.3.0$, and are run from the top directory of the ancillary files of this note, which are
self-contained: the scripts \nolinkurl{general/verify_general.py}, \nolinkurl{general/verify_pair.py} and
\nolinkurl{general/corollary_slack.py}, the parameter files in \nolinkurl{general/params/}, the summaries of the runs in
\nolinkurl{general/results/} and the logs in \nolinkurl{general/logs/}. The certificates \nolinkurl{zeta_K104_s13.json},
\nolinkurl{zeta_K80_s12_edge.json} and \nolinkurl{pair_Qi_X10.json} are used only here; the $K=128$ certificate
\nolinkurl{zeta_K128_s13.json} and the certificates for $\GR^2$ and $\GC$ (\nolinkurl{gammaR2_K80_s7.json},
\nolinkurl{pair_gammaR2_X240.json}, \nolinkurl{gammaC_K64_s9.json}, \nolinkurl{pair_gammaC_X40.json}) are also shipped, as identical
files, with the ancillary files of \cite{PaperI}. The script \nolinkurl{general/run_all.sh} repeats all of them (about
$25$ CPU-minutes), and \texttt{sha256sum -c SHA256SUMS} checks the inputs. The ancillary files are licensed under the Apache License 2.0 (copyright 2026 Nic Johns); the text is licensed under
CC BY 4.0.

\bibliographystyle{amsplain-doi}
\bibliography{refs}

\end{document}
